\documentclass[10pt,a4paper]{amsart}
\usepackage[margin=19mm]{geometry}
\usepackage{amsmath,amssymb,mathtools}
\usepackage[scr=rsfs,cal=euler]{mathalfa}
\usepackage{xfrac}
\usepackage[unicode,hidelinks,bookmarks=false]{hyperref}
\newcommand{\ie}{i.e.\ }
\newcommand{\cf}{cf.\ }
\newcommand{\resp}{respectively\ }
\newcommand{\Subsection}{\subsection}
\providecommand{\nobreakdash}{\nobreak\hbox{-}}
\setlength{\emergencystretch}{2em}
\multlinegap0pt
\usepackage{amssymb}
\newtheorem{definition}{Definition}
\newtheorem{theorem}{Theorem}
\newcommand{\interpret}[1]{\llbracket #1 \rrbracket}
\newcommand{\llbracket}{\left[\hspace{-0.15em}\left[}
\newcommand{\rrbracket}{\right]\hspace{-0.15em}\right]}
\title{Inexpressibility in Exp-Minus-Log}
\author{Mark Carney}
\date{Source: arXiv:2605.01636v1 (CC BY 4.0)}
\begin{document}
\maketitle

\noindent\emph{Source and licence.} Inexpressibility in Exp-Minus-Log. Version arXiv:2605.01636v1 (CC BY 4.0), distributed by arXiv under the Creative Commons Attribution 4.0 International licence, http://creativecommons.org/licenses/by/4.0/, as recorded in the arXiv licence metadata for this exact version and shown on the abstract page https://arxiv.org/abs/2605.01636v1. Full licence notice: \texttt{licenses/arxiv-2605.01636-CC-BY-4.0.txt}.\par\smallskip

\noindent\textit{Editorial scope.} Counted unit: the article's theorem thm:main (source0.tex lines 195--197). The article's complete original proof of it is delivered with it (source0.tex lines 198--200, one \texttt{proof} environment of 3 lines). No mathematics is written, repaired or re-derived: the statement and the proof are the article's own lines, in the article's own notation, and no retained line is substituted or re-flowed.  Retained uncounted as the source-local context the proof needs: lines 119--124; lines 126--128; lines 176--179.  Nothing was omitted from any retained span, so the package carries no omission record.  No retained line contains a citation, so the wrapper carries no bibliography and no bibliography entry.  No retained line contains a figure environment, an image-inclusion command or drawing code, so no image is delivered and the package declares no asset dependency.  Editorial formatting only, in addition: an AMS \texttt{amsart} preamble that restates the article's own declarations and macros used by the retained lines, namely the environments \texttt{definition}, \texttt{theorem} on separate counters, together with the standard packages those lines need.  The retained environments print in this document as Definition 1, Theorem 1 in order of appearance, whereas the article prints the same items with the numbering of its own full document.  The complete native source of the article as distributed in the arXiv e-print for this exact version is bundled unchanged as \texttt{sources/199649/source0.tex}.
\begin{definition}\label{def:omega}
    For any universal prefix-free Turing machine $U$, the halting probability
\begin{align}
    \Omega_U = \sum_{p\in \text{dom}(U)} 2^{-|p|} \in (0,1)
\end{align}
\end{definition}

\begin{theorem}[Chaitin, 1975]\label{thm:omega-non-comp}
    Let $U$ be any universal prefix-free Turing machine, and $\Omega_U$ defined per Def. \ref{def:omega}. $\Omega_U$ is a real number that is left-computably enumerable, but non-computable: no Turing machine, on input $k$, will output a rational within $2^{-k}$ of $\Omega_U$ for all $k$.
\end{theorem}

\begin{theorem}[Computability of EML values]
\label{thm:comp-eml-values}
    For every $\varphi \in \mathcal{E}$ with $\interpret{\varphi}$ defined, $\interpret{\varphi}$ is a computable complex number. In particular, $\mathrm{EML}_\mathbb{R} \subseteq \mathrm{Comp}_\mathbb{R}$.
\end{theorem}

\begin{theorem}\label{thm:main}
$$\Omega_U \notin EML_\mathbb{R}$$
\end{theorem}

\begin{proof}
    By Theorem \ref{thm:comp-eml-values}, $EML_\mathbb{R} \subseteq \mathtt{Comp}_\mathbb{R}$, and by Theorem \ref{thm:omega-non-comp}, $\Omega_U \in \mathbb{R} \setminus \mathtt{Comp}_\mathbb{R}$. Hence $\Omega_U \notin EML_\mathbb{R}$. Similarly, for any non-computable real $r$, $r\notin \mathtt{Comp}_\mathbb{R} \supseteq EML_\mathbb{R}$, therefore $\nexists \varphi \in \mathcal{E}$ such that $\interpret{\varphi} = r$.
\end{proof}

\end{document}
